\section{REM and Emotional \& Procedural Generalization}\label{sec:rem_generalization}

REM sleep serves a different but complementary function to NREM consolidation.

If NREM sleep saves the file, REM sleep opens it, edits it, and links it to other files. REM sleep is not primarily for consolidating raw facts, but for \textit{generalization} and \textit{integration}.

\begin{itemize}
    \item \textbf{Emotional Processing:} REM sleep is often called ``sleep to forget,~and sleep to remember''. During REM, the brain re-processes emotional memories but does so in a state of neurochemical suppression (low norepinephrine). This allows the brain to ``divorce the emotion from the memory'' \citep{yoo2007}. This is why, after a full night's sleep, a traumatic event can be recalled without re-experiencing the same raw, visceral emotional charge. It is a form of overnight therapy.
    
    \item \textbf{Procedural \& Creative Generalization:} REM sleep is crucial for ``understanding the rules'' of a new task, rather than just the facts. It helps us find hidden patterns and connections. This is why a person learning a new language (a procedural skill) or a musician practicing a new piece (a motor skill) shows dramatic improvement after REM sleep \citep{stickgold2013}. REM sleep moves beyond the literal replay of NREM and begins to ``remix'' information, forming novel associations. This is the neurobiological basis for why we ``sleep on a problem'' and wake up with a creative solution \citep{cai2009}.
\end{itemize}

These two functions—NREM consolidation and REM generalization—are the core of the brain's learning algorithm. This leads to the central thesis of this paper: that this algorithm is a solution to the ``overfitting'' problem, a concept we explore in~\ref{part:obh}.
